\section{Discussions}

\subsection{Sequence-Depth Duality}
\label{sec:depth-ttt}

Residual connections propagate information over depth via a fixed recurrence $\bm{h}_{l} = \bm{h}_{l-1} + f_{l-1}(\bm{h}_{l-1})$, much as RNNs propagate information over time.
Test-Time Training (TTT)~\cite{sun-2024-learning} formalizes the sequence side of this analogy (cf.\ Fast Weight Programmers~\cite{schmidhuber1992learning,munkhdalai-2019-metalearned}), casting each recurrent step as gradient descent on a self-supervised loss:
\begin{equation}
    \mathbf{W}_t = \mathbf{W}_{t-1} - \eta\,\nabla\ell(\mathbf{W}_{t-1};\, \bm{x}_t),
\end{equation}
where a slow network parameterizes $\ell$ and the state $\mathbf{W}$ is updated once per token.
When $f$ is linear, this reduces to vanilla linear attention $\mathbf{S}_{t} = \mathbf{S}_{t-1} + \bm{k}_t\bm{v}_t^\top$.
The standard residual exhibits the same additive form along depth, with $\bm{h}_l$ serving as the state and each layer $f_l$ acting as one ``gradient step.''

As noted by~\cite{chen2026postlayernorm}, this duality extends to richer variants (Table~\ref{tab:residual_compare}).
Data-dependent gates on the sequence side~\cite{sun-2023-retnet,yang-etal-2024-gla} correspond to Highway networks~\cite{srivastava2015highway} on the depth side;
the delta rule~\cite{schlag-2021-deltanet,yang-2025-gdn,zhang2025kimi} corresponds to DDL~\cite{zhang2026ddl};
and MRLA~\cite{fang2023cross} mirrors GLA's~\cite{yang-etal-2024-gla} gated linear attention.
These methods all refine the recurrent update while remaining within the recurrence paradigm.
AttnRes goes a step further and replaces depth-wise recurrence with direct cross-layer attention, just as Transformers replaced temporal recurrence with self-attention.
Since the number of layers in current architectures remains well within the practical regime of $\operatorname{softmax}$ attention, we adopt vanilla depth-wise attention. Incorporating more expressive yet memory-efficient (e.g. linear-complexity) alternatives is a natural direction for future work.

\begin{table*}[t]
  \centering
  \caption{
    Comparison of residual update mechanisms. \emph{Weight}: whether the mixing coefficients are architecture-fixed, learned-static (fixed after training), or input-dependent (dynamic). \emph{Source}: which earlier representations layer $l$ can access. Normalization is omitted from most formulas for clarity.
  }
  \label{tab:residual_compare}
  \small
  \setlength{\tabcolsep}{8pt}
  \renewcommand{\arraystretch}{1.3}
  \resizebox{\textwidth}{!}{%
    \begin{tabular}{l l l c c}
      \toprule
      \multicolumn{2}{l}{\textbf{Method}}                                   & \textbf{Update rule}                                                                                                                                                                         & \textbf{Weight}                                                                                                           & \textbf{Source}                                                                  \\
      \midrule
      \multicolumn{5}{@{}l}{\emph{Single-state recurrence: layer $l$ receives only $\bm{h}_{l-1}$}}                                                                                                                                                                                                                                                                                                                                                                                       \\
      \addlinespace[2pt]
      \multicolumn{2}{l}{Residual \citep{he2015resnet}}                     & $\bm{h}_{l} = \bm{h}_{l-1} + f_{l-1}(\bm{h}_{l-1})$                                                                                                                                          & Fixed                                                                                                                     & $\bm{h}_{l-1}$                                                                   \\
      \multicolumn{2}{l}{ReZero \citep{bachlechner2020rezero}}              & $\bm{h}_{l} = \bm{h}_{l-1} + \alpha_{l} \cdot f_{l-1}(\bm{h}_{l-1})$                                                                                                                         & Static                                                                                                                    & $\bm{h}_{l-1}$                                                                   \\
      \multicolumn{2}{l}{LayerScale \citep{touvron2021going}}               & $\bm{h}_{l} = \bm{h}_{l-1} + \mathrm{diag}(\bm{\lambda}_{l}) \cdot f_{l-1}(\bm{h}_{l-1})$                                                                                                    & Static                                                                                                                    & $\bm{h}_{l-1}$                                                                   \\
      \multicolumn{2}{l}{Highway \citep{srivastava2015highway}}             & $\bm{h}_{l} = (1{-}\bm{g}_{l}) \odot \bm{h}_{l-1} + \bm{g}_{l} \odot f_{l-1}(\bm{h}_{l-1})$                                                                                                  & Dynamic                                                                                                                   & $\bm{h}_{l-1}$                                                                   \\
      \multicolumn{2}{l}{DeepNorm \citep{wang2022deepnorm}}                 & $\bm{h}_{l} = \mathrm{Norm}(\alpha \, \bm{h}_{l-1} + f_{l-1}(\bm{h}_{l-1}))$                                                                                                                 & Fixed                                                                                                                     & $\bm{h}_{l-1}$                                                                   \\
      \multicolumn{2}{l}{KEEL \citep{chen2026postlayernorm}}                & $\bm{h}_{l} = \mathrm{Norm}(\alpha \, \bm{h}_{l-1} + f_{l-1}(\mathrm{Norm}(\bm{h}_{l-1})))$                                                                                                  & Fixed                                                                                                                     & $\bm{h}_{l-1}$                                                                   \\
      \midrule
      \multicolumn{5}{@{}l}{\emph{Multi-state recurrence: layer $l$ receives $m$ streams}}                                                                                                                                                                                                                                                                                                                                                                                                \\
      \addlinespace[2pt]
      \multicolumn{2}{l}{SiameseNorm \citep{li2026siamesenorm}}             & $\bm{h}^{1}_{l} {=} \mathrm{Norm}(\bm{h}^{1}_{l-1} {+} \bm{y}_{l-1})$;\; $\bm{h}^{2}_{l} {=} \bm{h}^{2}_{l-1} {+} \bm{y}_{l-1}$                                                              & Fixed                                                                                                                     & 2 streams                                                                        \\
      \multicolumn{2}{l}{HC/mHC \citep{zhu2025hyperconnections,xie2026mhc}} & $\mathbf{H}_{l} = \mathbf{H}_{l-1}\mathbf{A}_{l} + f_{l-1}(\mathbf{H}_{l-1}\bm{\alpha}_{l-1})\,\bm{\beta}_{l-1}^\top$                                                                        & Dynamic                                                                                                                   & $m$ streams                                                                      \\
      \multicolumn{2}{l}{DDL \citep{zhang2026ddl}}                          & $\mathbf{H}_{l} = (\mathbf{I} {-} \beta_{l} \bm{k}_{l}\bm{k}_{l}^\top)\mathbf{H}_{l-1} + \beta_{l} \bm{k}_{l}\bm{v}_{l}^\top$                                                                & Dynamic                                                                                                                   & $d_v$ streams                                                                    \\
      \midrule
      \multicolumn{5}{@{}l}{\emph{Cross-layer access: layer $l$ can access individual earlier-layer outputs}}                                                                                                                                                                                                                                                                                                                                                                             \\
      \addlinespace[2pt]
      \multicolumn{2}{l}{DenseNet \citep{huang2018densenet}}                & $\bm{h}_{l} = \mathrm{ConvPool}([\bm{h}_1;\; f_1(\bm{h}_1);\; \ldots;\; f_{l-1}(\bm{h}_{l-1})])$                                                                                             & Static                                                                                                                    & $[\bm{h}_1,\ldots,\bm{h}_{l-1}]$                                                 \\
      \multicolumn{2}{l}{DenseFormer \citep{pagliardini2024denseformer}}    & $\bm{h}_{l} = \alpha_{0 \to l}\,\bm{h}_1 + \sum_{i=1}^{l-1} \alpha_{i \to l}\,f_i(\bm{h}_i)$                                                                                                 & Static                                                                                                                    & $[\bm{h}_1,\ldots,\bm{h}_{l-1}]$                                                 \\
      \multicolumn{2}{l}{MRLA \citep{fang2023cross}$^{1}$}                  & $\bm{h}_{l} = \sum_{i=1}^{l-1} \sigma\!\bigl(\mathrm{ConvPool}(f_{l-1}(\bm{h}_{l-1}))\bigr)^{\!\top} \sigma\!\bigl(\mathrm{ConvPool}(f_{i}(\bm{h}_i))\bigr) \, \mathrm{Conv}(f_i(\bm{h}_i))$ & Dynamic                                                                                                                   & $[\bm{h}_1,\ldots,\bm{h}_{l-1}]$                                                 \\
      \midrule
      \rowcolor{gray!10}
                                                                            & Full$^{2}$                                                                                                                                                                                   & $\bm{h}_{l} \propto \sum_{i=0}^{l-1} \phi(\bm{w}_{l}, \bm{k}_i) \, \bm{v}_i$                                              & Dynamic                          & $[\bm{h}_1,\ldots,\bm{h}_{l-1}]$              \\
      \rowcolor{gray!10}
      \multirow{-2}{*}{AttnRes (ours)}                                      & Block$^{3}$                                                                                                                                                                                  & $\bm{h}_{l} \propto \sum_{i=0}^{n-1} \phi(\bm{w}_{l}, \bm{k}_i) \, \bm{v}_i + \phi(\bm{w}_{l}, \bm{k}_n^j) \, \bm{v}_n^j$ & Dynamic                          & $[\bm{b}_0,\ldots,\bm{b}_{n{-}1},\bm{b}_n^j]$ \\
      \bottomrule
    \end{tabular}%
  }% end resizebox
  \par\vspace{4pt}
  \raggedright\footnotesize
  $^{1}$\,ConvPool: pooling operation followed by convolution (channel projection).\\
  $^{2}$\,$\phi(\bm{q},\bm{k}) = \exp\left(\bm{q}^\top \operatorname{RMSNorm}(\bm{k})\right)$;\; $\bm{k}_i = \bm{v}_i$;\; $\bm{v}_0 = \bm{h}_1$,\; $\bm{v}_{i\geq1} = f_i(\bm{h}_i)$. $\operatorname{softmax}$ jointly normalized over all sources.\\
  $^{3}$\,Same $\phi$ and normalization as Full;\; $\bm{v}_i = \bm{b}_i$,\; $\bm{v}_n^j = \bm{b}_n^j$.
\end{table*}


\subsection{Residual Connections as Structured Matrices}
\label{sec:structured-depth}

The residual variants discussed above can all be viewed as weighted aggregations over previous layer outputs.
We formalize this with a \emph{depth mixing matrix} $\mathbf{M} \in \mathbb{R}^{L \times L}$, where $\mathbf{M}_{i \to l}$ is the weight that layer $l$ assigns to the output of layer $i$.
The variants differ in how these weights arise (fixed, learned, or input-dependent) and whether $\mathbf{M}$ is constrained to low rank or allowed to be dense.
The semiseparable rank of $\mathbf{M}$~\cite{mamba2} offers a unified lens for comparing them.

Concretely, the input to layer $l$ is $\bm{h}_l = \sum_{i=0}^{l-1} \mathbf{M}_{i \to l}\, \bm{v}_i$, where $\bm{v}_0 = \bm{h}_1$ (embedding) and $\bm{v}_i = f_i(\bm{h}_i)$ for $i \geq 1$.
Fig.~\ref{fig:depth-matrices} visualizes $\mathbf{M}$ for representative methods; we derive each below.
\begin{itemize}[leftmargin=*]
    \item Standard residual~\cite{he2015resnet}, $\bm{h}_{l} = \bm{h}_{l-1} + f_{l-1}(\bm{h}_{l-1})$. Expanding gives $\bm{h}_l = \sum_{i=0}^{l-1} \bm{v}_i$, so $\mathbf{M}_{i \to l} = 1$ for all $i < l$ and $\mathbf{M}$ is an all-ones lower-triangular matrix:
          \[
              \begin{bmatrix}
                  \bm{h}_1 \\ \bm{h}_2 \\ \vdots \\ \bm{h}_L
              \end{bmatrix}
              =
              \begin{bmatrix}
                  1      &        &        &   \\
                  1      & 1      &        &   \\
                  \vdots & \vdots & \ddots &   \\
                  1      & 1      & \cdots & 1
              \end{bmatrix}
              \begin{bmatrix}
                  \bm{v}_0 \\ \bm{v}_1 \\ \vdots \\ \bm{v}_{L-1}
              \end{bmatrix}
          \]
    \item Highway~\cite{srivastava2015highway}, $\bm{h}_{l} = (1{-}g_{l})\,\bm{h}_{l-1} + g_{l}\,f_{l-1}(\bm{h}_{l-1})$ (written here with scalar gates for clarity; the element-wise extension is straightforward). Defining the carry product $\gamma_{i \to l}^{\times} \coloneqq \prod_{j=i+1}^{l}(1-g_j)$, the weights are $\mathbf{M}_{0 \to l} = \gamma_{1 \to l}^{\times}$ for the embedding and $\mathbf{M}_{i \to l} = g_{i+1}\,\gamma_{i+1 \to l}^{\times}$ for $i \geq 1$. Since the cumulative products factor through scalar gates, $\mathbf{M}$ is 1-semiseparable~\cite{mamba2}, the same rank as the standard residual but with input-dependent weights. The weights sum to one by construction, making Highway a softmax-free depth-wise instance of stick-breaking attention~\cite{tan2025stickbreaking}.
    \item (m)HC~\cite{zhu2025hyperconnections,xie2026mhc} maintain $m$ parallel streams $\mathbf{H}_{l} \in \mathbb{R}^{d \times m}$, updated via
          \[
              \mathbf{H}_{l} = \mathbf{H}_{l-1} \mathbf{A}_{l} + f_{l-1}(\mathbf{H}_{l-1} \bm{\alpha}_{l-1})\, \bm{\beta}_{l-1}^\top,
          \]
          where $\mathbf{A}_{l} \in \mathbb{R}^{m \times m}$ is a learned transition matrix, $\bm{\alpha}_{l-1} \in \mathbb{R}^{m}$ mixes streams into a single input for $f_{l-1}$, and $\bm{\beta}_{l-1} \in \mathbb{R}^{m}$ distributes the output back across streams.
          Unrolling the recurrence gives the effective weight
          \begin{equation}
              \label{eq:hc-unroll}
              \mathbf{M}_{i \to l} = \bm{\beta}_{i}^\top \, \mathbf{A}_{i+1 \to l}^{\times} \, \bm{\alpha}_{l},
          \end{equation}
          where $\mathbf{A}_{i \to j}^{\times} \coloneqq \prod_{k=i+1}^{j} \mathbf{A}_k$.
          The $m \times m$ transitions render $\mathbf{M}$ $m$-semiseparable~\cite{mamba2}. mHC~\cite{xie2026mhc,yang2026mhclite} further constrains each $\mathbf{A}_l$ to be doubly stochastic, stabilizing the cumulative products across depth.
    \item Full AttnRes computes $\mathbf{M}_{i \to l} = \brickred{\alpha_{i \to l}}$ via $\phi(\bm{w}_l, \bm{k}_i) = \exp\left(\bm{w}_l^\top \operatorname{RMSNorm}(\bm{k}_i)\right)$ with normalization, where $\bm{k}_i = \bm{v}_i$ are input-dependent layer outputs, yielding a dense, rank-$L$ $\mathbf{M}$.
    \item Block AttnRes partitions layers into $N$ blocks $\mathcal{B}_1, \ldots, \mathcal{B}_N$. For sources $i$ in a completed earlier block $\mathcal{B}_n$, all share the block-level key/value $\bm{b}_n$, so $\mathbf{M}_{i \to l} = \brickred{\alpha_{n \to l}}$ for every $i \in \mathcal{B}_n$. Within the current block, each layer additionally attends over the evolving partial sum $\bm{b}_n^{i-1}$, introducing one extra distinct source per intra-block position. The effective rank of $\mathbf{M}$ therefore lies between $N$ and $N + S$ (where $S$ is the block size), interpolating between standard residual ($N{=}1$) and Full AttnRes ($N{=}L$).
\end{itemize}

\begin{figure*}[t]
    \centering
    \scriptsize

    \definecolor{CellBlue}{HTML}{E1EBF5}   % Soft slate blue
    \definecolor{CellSand}{HTML}{F5E8DF}   % Pale peach/orange
    \definecolor{CellPurple}{HTML}{EBE4F0} % Muted lilac
    \definecolor{CellMint}{HTML}{E1F0E5}   % Gentle sage green
    \definecolor{CellBlush}{HTML}{F5E1E6}  % Soft rose

    \renewcommand{\arraystretch}{1.4}
    \newlength{\gridw}\setlength{\gridw}{0.49\textwidth}
    \newlength{\cellw}\setlength{\cellw}{6.5em}
    \newcommand{\cell}[1]{\makebox[\cellw][c]{$#1$}}

    \newcommand{\roundcell}[3]{%
        \fill[#3, rounded corners=2pt] (#1-|#2) rectangle (\the\numexpr#1+1\relax-|\the\numexpr#2+1\relax);
    }

    \makebox[\gridw][c]{%
        \begin{tabular}{@{}c@{}}
            Highway \\[\medskipamount]
            $\begin{bNiceArray}{llll}[margin=2pt]
                     \cell{1}                         & \cell{}                               & \cell{}                               & \cell{}    \\
                     \cell{\gamma_{1 \to 2}^{\times}} & \cell{g_2}                            & \cell{}                               & \cell{}    \\
                     \cell{\gamma_{1 \to 3}^{\times}} & \cell{g_2\,\gamma_{2 \to 3}^{\times}} & \cell{g_3}                            & \cell{}    \\
                     \cell{\gamma_{1 \to 4}^{\times}} & \cell{g_2\,\gamma_{2 \to 4}^{\times}} & \cell{g_3\,\gamma_{3 \to 4}^{\times}} & \cell{g_4}
                 \end{bNiceArray}$
        \end{tabular}%
    }%
    \hfill
    \makebox[\gridw][c]{%
        \begin{tabular}{@{}c@{}}
            (m)HC \\[\medskipamount]
            $\begin{bNiceArray}{llll}[margin=2pt]
                     \cell{\bm{\beta}_0^\top \bm{\alpha}_1}                                 & \cell{}                                                                & \cell{}                                                              & \cell{}                                \\
                     \cell{\bm{\beta}_0^\top \mathbf{A}_{1 \to 2}^{\times} \bm{\alpha}_2}   & \cell{\bm{\beta}_1^\top \bm{\alpha}_2}                                 & \cell{}                                                              & \cell{}                                \\
                     \cell{\bm{\beta}_0^\top \mathbf{A}_{1 \to 3}^{\times}\, \bm{\alpha}_3} & \cell{\bm{\beta}_1^\top \mathbf{A}_{2 \to 3}^{\times} \bm{\alpha}_3}   & \cell{\bm{\beta}_2^\top \bm{\alpha}_3}                               & \cell{}                                \\
                     \cell{\bm{\beta}_0^\top \mathbf{A}_{1 \to 4}^{\times}\, \bm{\alpha}_4} & \cell{\bm{\beta}_1^\top \mathbf{A}_{2 \to 4}^{\times}\, \bm{\alpha}_4} & \cell{\bm{\beta}_2^\top \mathbf{A}_{3 \to 4}^{\times} \bm{\alpha}_4} & \cell{\bm{\beta}_3^\top \bm{\alpha}_4}
                 \end{bNiceArray}$
        \end{tabular}%
    }

    \vspace{14pt}

    \makebox[\gridw][c]{%
        \begin{tabular}{@{}c@{}}
            Full AttnRes \\[\medskipamount]
            $\begin{bNiceArray}{llll}[margin=2pt]
                     \CodeBefore
                     \begin{tikzpicture}
                        \roundcell{1}{1}{CellBlue}
                        \roundcell{2}{1}{CellBlue}
                        \roundcell{3}{1}{CellBlue}
                        \roundcell{4}{1}{CellBlue}
                        \roundcell{2}{2}{CellSand}
                        \roundcell{3}{2}{CellSand}
                        \roundcell{4}{2}{CellSand}
                        \roundcell{3}{3}{CellPurple}
                        \roundcell{4}{3}{CellPurple}
                        \roundcell{4}{4}{CellMint}
                    \end{tikzpicture}
                     \Body
                     \cell{\phi\left(\bm{w}_1, \bm{k}_0\right)} & \cell{}                              & \cell{}                            & \cell{}                            \\
                     \cell{\phi\left(\bm{w}_2, \bm{k}_0\right)} & \cell{\phi\left(\bm{w}_2, \bm{k}_1\right)} & \cell{}                            & \cell{}                            \\
                     \cell{\phi\left(\bm{w}_3, \bm{k}_0\right)} & \cell{\phi\left(\bm{w}_3, \bm{k}_1\right)} & \cell{\phi\left(\bm{w}_3, \bm{k}_2\right)} & \cell{}                            \\
                     \cell{\phi\left(\bm{w}_4, \bm{k}_0\right)} & \cell{\phi\left(\bm{w}_4, \bm{k}_1\right)} & \cell{\phi\left(\bm{w}_4, \bm{k}_2\right)} & \cell{\phi\left(\bm{w}_4, \bm{k}_3\right)}
                 \end{bNiceArray}$
        \end{tabular}%
    }%
    \hfill
    \makebox[\gridw][c]{%
        \begin{tabular}{@{}c@{}}
            Block AttnRes \\[\medskipamount]
            $\begin{bNiceArray}{llll}[margin=2pt]
                     \CodeBefore
                     \begin{tikzpicture}
                        \roundcell{1}{1}{CellBlue}
                        \roundcell{2}{1}{CellBlue}
                        \roundcell{3}{1}{CellBlue}
                        \roundcell{4}{1}{CellBlue}
                        \roundcell{2}{2}{CellSand}
                        \draw[dashed, rounded corners=2pt, line width=0.4pt] (2-|2) rectangle (3-|3);
                        \fill[CellBlush, rounded corners=2pt] (3-|2) rectangle (4-|4);
                        \fill[CellBlush, rounded corners=2pt] (4-|2) rectangle (5-|4);
                        \roundcell{4}{4}{CellMint}
                        \draw[dashed, rounded corners=2pt, line width=0.4pt] (4-|4) rectangle (5-|5);
                    \end{tikzpicture}
                     \Body
                     \cell{\phi\left(\bm{w}_1, \bm{k}_0\right)} & \cell{}                                  & \cell{}                            & \cell{}                                \\
                     \cell{\phi\left(\bm{w}_2, \bm{k}_0\right)} & \cell{\phi\left(\bm{w}_2, \bm{k}_1\right)} & \cell{}                            & \cell{}                                \\
                     \cell{\phi\left(\bm{w}_3, \bm{k}_0\right)} & \Block{1-2}{\cell{\phi\left(\bm{w}_3, \bm{k}_1+\bm{k}_2\right)}} &                            & \cell{}                                \\
                     \cell{\phi\left(\bm{w}_4, \bm{k}_0\right)} & \Block{1-2}{\cell{\phi\left(\bm{w}_4, \bm{k}_1 +\bm{k}_2\right)}} &                            & \cell{\phi\left(\bm{w}_4, \bm{k}_3\right)}
                 \end{bNiceArray}$
        \end{tabular}%
    }

    \caption{Depth mixing matrices $\mathbf{M}$ for four residual variants ($L{=}4$; Block AttnRes uses block size $S{=}2$). Highway is shown with scalar gates for clarity. AttnRes panels show unnormalized $\phi$ scores; background colors group entries that share the same source (Full AttnRes) or the same source block (Block AttnRes).}
    \label{fig:depth-matrices}
\end{figure*}

\paragraph{Practicality.}
The structured-matrix perspective serves two purposes.
First, it enables analytical insights that are not apparent from the recurrence form alone.
The input-dependent $\mathbf{M}$ of AttnRes, for instance, reveals depth-wise attention sinks (\S\ref{sec:analysis}), where certain layers consistently attract high weight regardless of input, mirroring the same phenomenon in sequence-wise attention~\cite{xiao2023efficient}.
Second, it informs new designs by exposing which properties of the kernel $\phi$ matter.
For example, when $\phi$ decomposes as $\phi(\bm{q}, \bm{k}) = \varphi(\bm{q})^\top \varphi(\bm{k})$ for some feature map $\varphi$~\cite{katharopoulos-2020-transformers}, depth-wise attention collapses into a recurrence---precisely the structure underlying the MRLA--GLA and DDL--DeltaNet correspondences noted above.

\paragraph{Prior Residuals as Depth-Wise Linear Attention}

The structured-matrix perspective further relates to the sequence-depth duality by showing that existing residual variants are, in effect, instances of \emph{linear} attention over the depth axis.
For example, the unrolled (m)HC weight $\mathbf{M}_{i \to l} = \boldsymbol{\beta}_i^\top \mathbf{A}_{i+1 \to l}^\times \boldsymbol{\alpha}_l$ (Eq.~\ref{eq:hc-unroll}) admits a natural attention interpretation in which $\boldsymbol{\alpha}_l$ plays the role of a query issued by layer $l$, $\boldsymbol{\beta}_i$ serves as a key summarizing the contribution of layer $i$, and the cumulative transition $\mathbf{A}_{i+1 \to l}^\times$ acts as a depth-relative positional operator~\cite{zhang2025kimi} governing the query--key interaction across intervening layers.
Notably, the $m$ parallel streams correspond to state expansion~\cite{qin-2024-hgrn2,mak2025residualmatrix} along the depth axis, expanding the recurrent state from $d$ to $d {\times}\, m$ and thereby increasing the semiseparable rank of $\mathbf{M}$.
\cite{xie2026mhcidentity} show that replacing $\mathbf{A}_{i+1 \to l}^{\times}$ with the identity matrix still yields competitive performance, highlighting the role of state expansion.
Through this lens, methods like (m)HC thus act as depth-wise \emph{linear} attention with matrix-valued states, while AttnRes acts as depth-wise $\operatorname{softmax}$ attention.





